\section{总结与延伸}

\subsection{核心讨论回顾}
\begin{table}[H]
\centering
\begin{tabular}{p{2.8cm}p{4.0cm}p{5.0cm}}
\toprule
\textbf{主题} & \textbf{panel 共识} & \textbf{代表性问题} \\
\midrule
产业背景 & 2025 是 Agent 从概念走向闭环的一年 & 为什么资本、产品、研究同时起量？ \\
模型训练 & 要把模型做成 native agentic，而非只靠 prompting & 哪些能力该内化进模型参数？ \\
系统工程 & LM system 仍然关键，模型越强越需要外层系统承接 & memory、workflow、权限如何设计？ \\
Benchmark / RL & solid benchmark 和 environment 决定方法是否能稳定比较 & RL 的 reward 与环境如何构造？ \\
未来方向 & 从固定模板走向自主探索，继续 scale acting 与 environment & 2026 年最重要的突破口在哪里？ \\
\bottomrule
\end{tabular}
\caption{把整场讨论压缩成五个判断轴}
\end{table}

\subsection{值得带走的三条结论}
\begin{enumerate}
  \item Agent 的竞争不再是单一模型能力竞争，而是模型能力、系统接口与环境构造的联动竞争。
  \item 2025 年最重要的研究变化，是开始认真把 agentic 能力当成训练对象，而不只是 prompt 产物。
  \item 真正限制 Agent 规模化落地的，往往不是 demo 是否惊艳，而是 environment、benchmark、memory、权限与部署是否足够扎实。
\end{enumerate}

\begin{importantbox}{这场 panel 最有价值的地方}
它没有陷入 ``哪个框架更好'' 的表层争论，而是把 Agent 当成一个完整系统来讨论：模型、工具、环境、训练、产品与行业落地都必须一起看。只有这样，``Agent'' 才不只是一个流行名词。
\end{importantbox}

\subsection{拓展阅读}
\begin{itemize}
  \item Deep Research / Deep Analysis 一类数据分析 Agent 系统公开资料
  \item Mobile-Agent、GUI Agent、Computer Use 相关论文与评测
  \item TinyZero、APR、multi-agent collaboration 与 Agentic training 方向论文
  \item 关于 world model / environment model 用于 Agent 训练与评测的近期工作
  \item 代码 Agent 产品实践：Cursor、Claude Code 等系统化工作流设计资料
\end{itemize}

\end{document}
